% ═══════════════════════════════════════════════════════════════
% MT-07a-ML: Musterlösung Minitest Lebensmittel (Comida)
% Spanisch Klasse 7 - Sequenz 7a
% ═══════════════════════════════════════════════════════════════
% Dauer: 10-15 Minuten
% Punkte: 25 BE
% Inhalt: Lebensmittel, Artikel, Mengenangaben, Kategorisierung
% ═══════════════════════════════════════════════════════════════

\documentclass[11pt, a4paper]{article}

\newif\ifswmodus
\swmodusfalse

\usepackage[utf8]{inputenc}
\usepackage[T1]{fontenc}
\usepackage[ngerman]{babel}
\usepackage{helvet}
\renewcommand{\familydefault}{\sfdefault}

\usepackage[margin=2cm]{geometry}
\usepackage{enumitem}
\usepackage{tabularx}
\usepackage{booktabs}

\usepackage{xcolor}
\usepackage{tcolorbox}
\tcbuselibrary{skins}

\usepackage{fancyhdr}
\usepackage{lastpage}
\usepackage{needspace}

% FARBEN
\definecolor{spanisch-color}{HTML}{c41e3a}
\definecolor{grau-color}{HTML}{888888}
\definecolor{hellgrau-color}{HTML}{f5f5f5}
\definecolor{orange-color}{HTML}{b35c00}
\definecolor{loesung-color}{HTML}{c62828}

\definecolor{sw-dark}{HTML}{000000}
\definecolor{sw-gray}{HTML}{666666}
\definecolor{sw-white}{HTML}{ffffff}

\ifswmodus
    \colorlet{spanisch}{sw-dark}
    \colorlet{grau}{sw-gray}
    \colorlet{hellgrau}{sw-white}
    \colorlet{orange}{sw-dark}
    \colorlet{loesung}{sw-dark}
\else
    \colorlet{spanisch}{spanisch-color}
    \colorlet{grau}{grau-color}
    \colorlet{hellgrau}{hellgrau-color}
    \colorlet{orange}{orange-color}
    \colorlet{loesung}{loesung-color}
\fi

\newcommand{\fachfarbe}{spanisch}

\newcommand{\thema}{Minitest: Comida -- MUSTERLÖSUNG}
\newcommand{\fach}{Spanisch}
\newcommand{\klasse}{7}
\newcommand{\maxpunkte}{25}
\newcommand{\zeit}{10-15 Minuten}
\newcommand{\datum}{\today}

\pagestyle{fancy}
\fancyhf{}
\renewcommand{\headrulewidth}{0.4pt}
\renewcommand{\footrulewidth}{0.4pt}

\lhead{\fach{} -- Klasse \klasse}
\chead{\textbf{MT-07a-ML}}
\rhead{\datum}

\lfoot{\zeit}
\cfoot{}
\rfoot{Seite \thepage\ von \pageref{LastPage}}

\newcommand{\punktebox}[1]{\hfill\fbox{\small \_/#1}}
\newcommand{\luecke}[1]{\rule{#1}{0.4pt}}
\newcommand{\lsg}[1]{\textcolor{loesung}{\textbf{#1}}}

\newtcolorbox{infobox}{
    colback=hellgrau,
    colframe=\fachfarbe,
    boxrule=1pt,
    arc=0pt,
    left=8pt, right=8pt, top=8pt, bottom=8pt
}

\newtcolorbox{hinweis}{
    colback=white,
    colframe=orange,
    boxrule=1pt,
    arc=0pt,
    left=5pt, right=5pt, top=5pt, bottom=5pt
}

\newenvironment{aufgabe}[1]{%
    \needspace{5\baselineskip}%
    \nopagebreak[4]%
    \vspace{10pt}
    \noindent\textbf{\textcolor{\fachfarbe}{Aufgabe #1}}
    \nopagebreak[4]%
    \vspace{5pt}
    \nopagebreak[4]%
    \begin{list}{}{\leftmargin=0pt}
    \item[]
}{%
    \end{list}
}

\newcommand{\es}[1]{\textcolor{\fachfarbe}{\textbf{#1}}}

\begin{document}

% ═══════════════════════════════════════════════════════════════
% KOPFBEREICH
% ═══════════════════════════════════════════════════════════════

\begin{center}
    {\Large\bfseries\textcolor{loesung}{\thema}}
\end{center}

\vspace{5pt}

\noindent
\begin{tabularx}{\textwidth}{|X|X|X|}
\hline
\textbf{Name:} \lsg{Musterlösung} &
\textbf{Klasse:} \klasse &
\textbf{Datum:} \luecke{2cm} \\
\hline
\end{tabularx}

\vspace{10pt}

\begin{infobox}
\textbf{Zeit:} \zeit{} | \textbf{Punkte:} \maxpunkte{} BE | \textbf{Hilfsmittel:} keine
\end{infobox}

\vspace{15pt}

% ═══════════════════════════════════════════════════════════════
% AUFGABE 1: Zuordnung Lebensmittel (6 BE)
% ═══════════════════════════════════════════════════════════════

\begin{aufgabe}{1 -- Lebensmittel zuordnen}
Verbinde das deutsche Wort mit der spanischen Übersetzung. \punktebox{6}
\end{aufgabe}

\begin{center}
\begin{tabular}{rcl}
1. Apfel & \lsg{--- 2 ---} & el plátano \\[0.8em]
2. Tomate & \lsg{--- 1 ---} & la manzana \\[0.8em]
3. Orange & \lsg{--- 6 ---} & la leche \\[0.8em]
4. Milch & \lsg{--- 3 ---} & la lechuga \\[0.8em]
5. Banane & \lsg{--- 5 ---} & el tomate \\[0.8em]
6. Salat & \lsg{--- 4 ---} & la naranja \\
\end{tabular}
\end{center}

\textit{Alternative Notation:}
\begin{itemize}[nosep]
    \item 1. Apfel = \lsg{la manzana}
    \item 2. Tomate = \lsg{el tomate}
    \item 3. Orange = \lsg{la naranja}
    \item 4. Milch = \lsg{la leche}
    \item 5. Banane = \lsg{el plátano}
    \item 6. Salat = \lsg{la lechuga}
\end{itemize}

\vspace{15pt}

% ═══════════════════════════════════════════════════════════════
% AUFGABE 2: Artikel einsetzen (6 BE)
% ═══════════════════════════════════════════════════════════════

\begin{aufgabe}{2 -- Artikel einsetzen}
Ergänze den richtigen Artikel: \es{el} oder \es{la}. \punktebox{6}
\end{aufgabe}

\noindent
a) \lsg{la} manzana \hfill (Apfel)

\vspace{0.8em}

\noindent
b) \lsg{el} tomate \hfill (Tomate)

\vspace{0.8em}

\noindent
c) \lsg{el} pollo \hfill (Hähnchen)

\vspace{0.8em}

\noindent
d) \lsg{la} leche \hfill (Milch)

\vspace{0.8em}

\noindent
e) \lsg{el} queso \hfill (Käse)

\vspace{0.8em}

\noindent
f) \lsg{la} zanahoria \hfill (Karotte)

\vspace{15pt}

% ═══════════════════════════════════════════════════════════════
% AUFGABE 3: Mengenangaben (5 BE)
% ═══════════════════════════════════════════════════════════════

\begin{aufgabe}{3 -- Mengenangaben}
Wähle die passende Mengenangabe für jedes Lebensmittel. \punktebox{5}
\end{aufgabe}

\begin{hinweis}
\textbf{Mengenangaben:} un kilo de | medio kilo de | un litro de | una botella de | un paquete de
\end{hinweis}

\vspace{0.5em}

\noindent
a) \lsg{un kilo de / medio kilo de} manzanas

\vspace{0.8em}

\noindent
b) \lsg{un litro de} leche

\vspace{0.8em}

\noindent
c) \lsg{una botella de / un litro de} agua

\vspace{0.8em}

\noindent
d) \lsg{un paquete de} arroz

\vspace{0.8em}

\noindent
e) \lsg{un kilo de / medio kilo de} tomates

\vspace{15pt}

% ═══════════════════════════════════════════════════════════════
% AUFGABE 4: Kategorisierung (8 BE)
% ═══════════════════════════════════════════════════════════════

\begin{aufgabe}{4 -- Kategorien zuordnen}
Ordne die Lebensmittel in die richtigen Kategorien ein. \punktebox{8}
\end{aufgabe}

\begin{hinweis}
\textbf{Wörter:} el pollo, la naranja, el queso, la zanahoria, el pan, el pescado, el plátano, la cebolla
\end{hinweis}

\vspace{0.5em}

\noindent
\begin{tabularx}{\textwidth}{|X|X|}
\hline
\textbf{La fruta} (Obst) & \textbf{La verdura} (Gemüse) \\
\lsg{la naranja, el plátano} & \lsg{la zanahoria, la cebolla} \\[1em]
\hline
\textbf{La carne / El pescado} & \textbf{Otros} (Andere) \\
\lsg{el pollo, el pescado} & \lsg{el queso, el pan} \\[1em]
\hline
\end{tabularx}

\vspace{0.5em}

\textit{Bewertung: Je 2 richtige Wörter pro Kategorie = 2 BE (insgesamt 8 BE)}

% ═══════════════════════════════════════════════════════════════
% PUNKTEÜBERSICHT
% ═══════════════════════════════════════════════════════════════

\vspace{25pt}
\hrule
\vspace{10pt}

\noindent
\begin{tabularx}{\textwidth}{|X|c|c|c|c||c|}
\hline
\textbf{Aufgabe} & 1 & 2 & 3 & 4 & \textbf{Gesamt} \\
\hline
\textbf{max. BE} & 6 & 6 & 5 & 8 & \textbf{25} \\
\hline
\textbf{erreicht} & & & & & \\[0.8em]
\hline
\end{tabularx}

\vspace{15pt}

% ═══════════════════════════════════════════════════════════════
% 14-NP-SKALA
% ═══════════════════════════════════════════════════════════════

\noindent\textbf{14-NP-Bewertungsskala (Sek I):}

\vspace{0.5em}

\noindent
\begin{tabularx}{\textwidth}{|c|c|c|c|c|c|c|c|c|c|c|c|c|c|c|}
\hline
\textbf{NP} & 14 & 13 & 12 & 11 & 10 & 9 & 8 & 7 & 6 & 5 & 4 & 3 & 2 & 1 \\
\hline
\textbf{BE} & 25 & 24 & 23 & 22 & 20 & 18 & 17 & 15 & 13 & 12 & 10 & 8 & 7 & 5 \\
\hline
\end{tabularx}

\vspace{10pt}

\begin{flushright}
\textbf{Note:} \luecke{2cm}
\end{flushright}

\end{document}
